\قسمت{تئوری کلاه}

از شما خواسته شده است یک برنامه موبایل توسعه دهید که در آن با یک سرور در ارتباط هستید. این برنامه به کاربران پست‌های مختلفی را نشان داده و کاربران می‌توانند آن‌ها لایک کنند.
برنامه شما به این شیوه کار می‌کند که اگر ارتباط با سرور امکان پذیر نباشد تغییرات شما را ذخیره کرده و در اولین اتصال صحیح آن را در اختیار کاربران قرار می‌دهد.
مدیر شما سفارش جدی کرده است که کاربران قطعی سرور را احساس نکنند.
شکایات مردمی حاکی از این می‌باشد که کاربران از رفتار عجیب اپلیکیشن و تغییرات ناگهانی تعداد لایک‌ها، ناراضی هستند.
آیا می‌توانید نظر مدیر و کاربران را همزمان جلب کنید؟ این بحث را با تئوری \متن‌لاتین{CAP} که در درس توضیح داده شد بررسی کنید و این تئوری را نیز مختصرا توضیح دهید.

\نمره{۳}

% پاسخ پیشنهادی است و پیش از استفاده در تصحیح نیاز به بازبینی دارد.
\begin{پاسخ}

تئوری \متن‌لاتین{CAP} می‌گوید یک سامانه‌ی توزیع‌شده در زمان بروز افراز شبکه\پانویس{Partition} تنها می‌تواند یکی از سازگاری\پانویس{Consistency} یا در دسترس بودن\پانویس{Availability} را نگه دارد.

طراحی گفته‌شده \متن‌لاتین{A} را انتخاب کرده است: در نبود شبکه هم کاربر می‌تواند لایک کند و تغییرات بعدا ارسال می‌شوند.
پرش ناگهانی تعداد لایک‌ها دقیقا هزینه‌ی همین انتخاب است و نمی‌توان همزمان هر دو خواسته را به شکل کامل برآورده کرد.

آنچه می‌توان بهبود داد تجربه‌ی کاربری است، نه خود قضیه: نمایش صریح وضعیت «در انتظار همگام‌سازی»، ادغام حساب‌شده‌ی تغییرات (شمارنده‌های تجمیع‌پذیر) و پرهیز از نمایش عددی که بعدا کاهش پیدا می‌کند.

\شروع{فقرات}
\فقره توضیح مختصر \متن‌لاتین{CAP}: یک سوم نمره
\فقره تشخیص اینکه طراحی فعلی \متن‌لاتین{AP} است: یک سوم نمره
\فقره بیان اینکه جلب همزمان هر دو ممکن نیست و تنها تجربه‌ی کاربری قابل بهبود است: یک سوم نمره
\پایان{فقرات}

\end{پاسخ}
